\documentclass[11pt]{article}

% ============================================================
% Packages
% ============================================================
\usepackage[utf8]{inputenc}
\usepackage[T1]{fontenc}
\usepackage{lmodern}
\usepackage{amsmath,amssymb,amsthm,mathtools}
\usepackage{enumitem}
\usepackage{geometry}
\usepackage{microtype}
\usepackage{hyperref}
\usepackage{xcolor}

\geometry{margin=1in}

% ============================================================
% Commands
% ============================================================
\newcommand{\R}{\mathbb{R}}
\newcommand{\N}{\mathbb{N}}
\newcommand{\dd}{\,dx}
\newcommand{\norm}[1]{\left\|#1\right\|}
\newcommand{\abs}[1]{\left|#1\right|}
\newcommand{\pc}{p_{\mathrm c}}

% ============================================================
% Theorem environments
% ============================================================
\newtheorem{theorem}{Theorem}[section]
\newtheorem{proposition}[theorem]{Proposition}
\newtheorem{lemma}[theorem]{Lemma}
\newtheorem{corollary}[theorem]{Corollary}

\theoremstyle{definition}
\newtheorem{remark}[theorem]{Remark}

% ============================================================
% Title and authors
% ============================================================

\title{On the Ambrosetti--Brezis--Cerami conjecture}

\author{
Anderson L.A.~de Araujo\thanks{
Universidade Federal de Vi\c{c}osa,
Departamento de Matem\'atica,
Avenida Peter Henry Rolfs,
Vi\c{c}osa, MG, Brazil, CEP 36570-000.
E-mail: \texttt{anderson.araujo@ufv.br}}
\and
Luiz F.O.~Faria\thanks{
Universidade Federal de Juiz de Fora,
Departamento de Matem\'atica,
Rua Jos\'e Louren\c{c}o Kelmer, s/n,
Juiz de Fora, MG, Brazil, CEP 36036-900.
E-mail: \texttt{luiz.faria@ufjf.br}}
}

\date{}

% ============================================================
\begin{document}

\maketitle

\begin{abstract}
We revisit the small-parameter uniqueness problem for the
concave--convex Dirichlet equation
\[
-\Delta u=\lambda u^q+u^p \quad\text{in }B_1,
\qquad u>0 \quad\text{in }B_1,
\qquad u=0 \quad\text{on }\partial B_1,
\]
where $B_1\subset\mathbb{R}^N$, $N\geq3$, $0<q<1$, and
$p>(N+2)/(N-2)$. This problem originates from an open question of
Ambrosetti, Brezis and Cerami concerning multiplicity in the
Sobolev-supercritical regime, for which small-parameter uniqueness
was subsequently stated by Zhao and Zhong through an approach based
on radial reduction and singular solutions. We give a new,
self-contained proof based on the radial first-zero map. Small
amplitudes are controlled by the pure concave limit, intermediate
amplitudes by a uniform lower bound for the first zero, and large
amplitudes by an Emden--Fowler analysis of the supercritical
Lane--Emden equation. These estimates yield a unique positive radial
solution for all sufficiently small $\lambda>0$. A symmetry argument
adapted to the non-Lipschitz concave term then gives uniqueness among
all positive classical solutions. Moreover,
\[
\|u_\lambda\|_{L^\infty(B_1)}
\sim
\rho_q^{-\frac{2}{1-q}}\lambda^{\frac1{1-q}}
\qquad\text{as }\lambda\to0^+,
\]
where $\rho_q$ is the first positive zero of the normalized pure
concave profile.
\end{abstract}

%%%%%%%%%%%%%%%%%%%%%%%%%%%%%%%%%%%%%%%%%%%%%%%%%%%%%%%%%%%%%%%%%
%%%%%%%%%%%%%%%%%%%%%%%%%%%%%%%%%%%%%%%%%%%%%%%%%%%%%%%%%%%%%%%%%
\section{Introduction and main result}
%%%%%%%%%%%%%%%%%%%%%%%%%%%%%%%%%%%%%%%%%%%%%%%%%%%%%%%%%%%%%%%%%
Let
\[
B_1:=B_1(0)\subset\mathbb{R}^N,
\qquad N\geq3,
\]
and consider the concave--convex Dirichlet problem
\begin{equation}
\label{eq:ABC}
\begin{cases}
-\Delta u=\lambda u^q+u^p,
&\text{in }B_1,\\
u>0,
&\text{in }B_1,\\
u=0,
&\text{on }\partial B_1,
\end{cases}
\end{equation}
where
\begin{equation}
\label{eq:hyp}
0<q<1,
\qquad
p>\pc:=\frac{N+2}{N-2},
\qquad
\lambda>0.
\end{equation}

The study of~\eqref{eq:ABC} goes back to the classical work of
Ambrosetti, Brezis and Cerami~\cite{ABC}, where the combined effect
of a concave and a convex nonlinearity was shown to produce a rich
existence and multiplicity theory. In particular, they proved the
existence of a minimal positive solution for sufficiently small
$\lambda>0$ and obtained a second positive solution in the
Sobolev-subcritical and critical range
\[
1<p\leq \frac{N+2}{N-2}.
\]
Their analysis left open the corresponding multiplicity question in
the Sobolev-supercritical regime: if
\[
p>\frac{N+2}{N-2}
\]
and the domain is a ball, does a second positive solution still exist
for sufficiently small $\lambda>0$?

The supercritical case is substantially different from the
subcritical and critical ones. The usual compact variational
framework is no longer available, and the possible behavior of
large-amplitude solutions becomes an essential part of the problem.
A radial multiplicity result was subsequently stated by
Ma~\cite{Ma1996} for $0<q<1<p$, with a formulation including the
supercritical range. On the other hand, Zhao and
Zhong~\cite{ZhaoZhong2002} later reached the opposite conclusion for
the problem considered here: for
\[
0<q<1,
\qquad
p>\frac{N+2}{N-2},
\]
they stated uniqueness of the positive solution for sufficiently
small $\lambda>0$, thus giving a negative answer to the open question
of Ambrosetti, Brezis and Cerami.

The argument of Zhao--Zhong is based on the asymptotic behavior of
large solutions and on a distinguished singular radial solution.
After reducing the problem to the radial equation, they associate
with the singular solution a parameter $\lambda_*>0$ and assert that
every sequence of regular solutions $(\lambda_n,u_n)$ satisfying
\[
\|u_n\|_{L^\infty(B_1)}\longrightarrow\infty
\]
must satisfy
\[
\lambda_n\longrightarrow\lambda_*>0,
\]
with $u_n$ converging to the singular profile away from the origin.
Their small-parameter uniqueness result is then obtained by combining
this asymptotic statement with the assertion that any solution
different from the minimal one must satisfy
\[
\|v_\lambda\|_{L^\infty(B_1)}
\longrightarrow\infty
\qquad\text{as }\lambda\to0^+.
\]

Some steps in the published proof of this argument are presented only
in abbreviated form. In particular, the latter blow-up alternative,
which is essential in the deduction of small-parameter uniqueness,
is stated in~\cite{ZhaoZhong2002} without proof, with reference to an
argument from the earlier Ambrosetti--Brezis--Cerami theory. Moreover,
the asymptotic identification of an arbitrary unbounded sequence with
the distinguished singular orbit involves limiting steps that are
only briefly justified in the published proof. We do not claim that
these points cannot be completed, and the purpose of the present work
is not to provide a corrigendum to~\cite{ZhaoZhong2002}. Instead, we
develop a different, self-contained approach which avoids the
singular-solution limiting argument altogether.

The main idea is to study directly the geometry of the first-zero map
associated with the radial shooting problem. After an appropriate
scaling, problem~\eqref{eq:ABC} is transformed into
\begin{equation}
\label{eq:scaled-intro}
\begin{cases}
-\Delta U=U^q+U^p,
&\text{in }B_R,\\
U>0,
&\text{in }B_R,\\
U=0,
&\text{on }\partial B_R,
\end{cases}
\end{equation}
where
\begin{equation}
\label{eq:Rlambda-intro}
R=\lambda^{\frac{p-1}{2(p-q)}}.
\end{equation}
Thus the limit $\lambda\to0^+$ corresponds precisely to the
small-radius limit $R\to0^+$.

For $A>0$, let $U_A$ denote the radial solution of
\begin{equation}
\label{eq:shooting-intro}
U_A''
+\frac{N-1}{r}U_A'
+U_A^q+U_A^p=0,
\qquad
U_A(0)=A,
\qquad
U_A'(0)=0,
\end{equation}
and let $R(A)$ be its first positive zero. Positive radial solutions
of~\eqref{eq:scaled-intro} are therefore characterized by
\[
R(A)=R.
\]
The uniqueness problem is reduced to understanding the first-zero map
$A\mapsto R(A)$ when $R$ is small.

Our analysis separates naturally into three amplitude regimes. For
small $A$, the concave term dominates. With the rescaling
\[
U_A(r)
=
A\,V_A\!\left(A^{-\frac{1-q}{2}}r\right),
\]
the functions $V_A$ converge to the normalized solution $V$ of
\begin{equation}
\label{eq:pure-concave-intro}
V''
+\frac{N-1}{r}V'
+V^q=0,
\qquad
V(0)=1,
\qquad
V'(0)=0.
\end{equation}
If $\rho_q$ denotes the first positive zero of $V$, we obtain
\begin{equation}
\label{eq:R-small-intro}
R(A)
=
\rho_q A^{\frac{1-q}{2}}(1+o(1))
\qquad\text{as }A\to0^+.
\end{equation}
In addition, a Picone-type argument gives uniqueness below the
explicit amplitude
\begin{equation}
\label{eq:Astar-intro}
A_*=
\left(
\frac{1-q}{p-1}
\right)^{\frac1{p-q}},
\end{equation}
which is exactly the interval on which
\[
s\longmapsto
\frac{s^q+s^p}{s}
=
s^{q-1}+s^{p-1}
\]
is strictly decreasing.

The essential issue is therefore to rule out additional solutions
whose central values stay away from zero. On every fixed compact
amplitude interval
\[
0<A_0\leq A\leq A_1<\infty,
\]
an elementary shooting estimate provides a uniform positive lower
bound for $R(A)$. The genuinely supercritical mechanism appears in
the large-amplitude regime. Setting
\[
U_A(r)
=
A\,v_\varepsilon
\left(A^{\frac{p-1}{2}}r\right),
\qquad
\varepsilon=A^{q-p},
\]
leads to
\[
v_\varepsilon''
+\frac{N-1}{r}v_\varepsilon'
+v_\varepsilon^p
+\varepsilon v_\varepsilon^q
=0.
\]
As $A\to\infty$, equivalently $\varepsilon\to0^+$, this equation
approaches the supercritical Lane--Emden equation
\[
v''
+\frac{N-1}{r}v'
+v^p=0.
\]

The Emden--Fowler transformation converts the limiting equation into
an autonomous planar system possessing a positive equilibrium. The
condition
\[
p>\frac{N+2}{N-2}
\]
is precisely what makes the damping coefficient positive. The
equilibrium is consequently asymptotically stable, and a robust
trapping argument shows that the small non-autonomous perturbation
generated by the concave term cannot drive the corresponding orbit
to zero before a sufficiently long Emden--Fowler time. After
returning to the original variables, the scaling factors cancel
exactly and yield a positive lower bound for the first zero of all
large-amplitude profiles.

Combining the intermediate- and large-amplitude estimates, we prove
the fundamental separation property
\begin{equation}
\label{eq:separation-intro}
\inf_{A\geq A_0}R(A)>0
\qquad
\text{for every }A_0>0.
\end{equation}
Hence, if $R>0$ is sufficiently small, every positive radial solution
of~\eqref{eq:scaled-intro} necessarily belongs to the small-amplitude
branch, where uniqueness is already available. In particular, the
argument excludes a second solution directly at the level of the
shooting diagram; no passage to a singular limiting solution and no
a priori blow-up alternative for a hypothetical second branch are
required.

This gives a different dynamical explanation of the small-parameter
phenomenon in the supercritical regime. The contribution of the
present work is therefore not only a new proof of the uniqueness
statement. The first-zero analysis provides a self-contained
mechanism that controls the complete amplitude range, makes explicit
where the supercritical assumption enters the argument, and yields
quantitative information on the unique small branch. In particular,
combining~\eqref{eq:Rlambda-intro} and~\eqref{eq:R-small-intro}, we
obtain
\begin{equation}
\label{eq:asymptotic-intro}
\|u_\lambda\|_{L^\infty(B_1)}
\sim
\rho_q^{-\frac{2}{1-q}}
\lambda^{\frac1{1-q}}
\qquad\text{as }\lambda\to0^+.
\end{equation}

There remains one point that requires some care. Since $0<q<1$, the
reaction
\[
f_\lambda(s)=\lambda s^q+s^p
\]
is not locally Lipschitz at $s=0$. Thus the passage from radial
uniqueness to uniqueness among all positive solutions should not be
based without further justification on a moving-plane theorem whose
standard formulation assumes local Lipschitz continuity of the
reaction. We use instead a symmetry result applicable to the present
structure, observing that
\[
f_\lambda(s)
=
(s^+)^p+\lambda(s^+)^q
\]
is the sum of a locally Lipschitz function and a nondecreasing
function. It follows that every positive classical solution of
\eqref{eq:ABC} in the ball is radial and strictly decreasing in the
radial variable. Consequently, the radial uniqueness obtained from
the shooting analysis yields uniqueness among all positive classical
solutions.

Our main result can now be stated as follows.

\begin{theorem}
\label{thm:main}
Assume~\eqref{eq:hyp}. Then there exists $\lambda_*>0$ such that,
for every
\[
0<\lambda<\lambda_*,
\]
problem~\eqref{eq:ABC} admits exactly one positive classical solution,
denoted by $u_\lambda$. Moreover, $u_\lambda$ is radial and, writing
\[
u_\lambda(x)=U_\lambda(|x|),
\]
one has
\[
U_\lambda'(r)<0
\qquad
\text{for every }0<r<1.
\]
Furthermore,
\[
\|u_\lambda\|_{L^\infty(B_1)}
\sim
\rho_q^{-\frac{2}{1-q}}
\lambda^{\frac1{1-q}}
\qquad
\text{as }\lambda\to0^+,
\]
where $\rho_q$ denotes the first positive zero of the solution of
\[
V''
+\frac{N-1}{r}V'
+V^q=0,
\qquad
V(0)=1,
\qquad
V'(0)=0.
\]
\end{theorem}

The proof of Theorem~\ref{thm:main} is organized around four main
ingredients. First, we perform the scaling that transforms the
small-parameter problem into the parameter-free equation on a small
ball. Second, we analyze the pure concave limit and establish
uniqueness below the explicit amplitude threshold~\eqref{eq:Astar-intro}.
Third, we study the first-zero map over the intermediate and
large-amplitude regimes, the latter through the Emden--Fowler
dynamics of the supercritical Lane--Emden equation. Finally, we apply
the appropriate symmetry principle for the non-Lipschitz
concave--convex reaction to pass from radial uniqueness to full
uniqueness.


%%%%%%%%%%%%%%%%%%%%%%%%%%%%%%%%%%%%%%%%%%%%%%%%%%%%%%%%%%%%%%%%%
\section{Scaling and the shooting formulation}
%%%%%%%%%%%%%%%%%%%%%%%%%%%%%%%%%%%%%%%%%%%%%%%%%%%%%%%%%%%%%%%%%

Set
\begin{equation}
\label{eq:gamma}
\gamma
:=
\frac{p-1}{2(p-q)}.
\end{equation}
For $\lambda>0$, consider the change of variables
\begin{equation}
\label{eq:scaling}
u(x)
=
\lambda^{\frac1{p-q}}
U\!\left(\lambda^\gamma x\right).
\end{equation}
Then, under this change of variables, $u$ solves~\eqref{eq:ABC}
if and only if $U$ solves
\begin{equation}
\label{eq:scaled}
\begin{cases}
-\Delta U=U^q+U^p,
&\text{in }B_R,\\
U>0,
&\text{in }B_R,\\
U=0,
&\text{on }\partial B_R,
\end{cases}
\end{equation}
where
\begin{equation}
\label{eq:Rlambda}
R=\lambda^\gamma.
\end{equation}

Indeed,
\[
-\Delta_xu
=
\lambda^{\frac1{p-q}+2\gamma}(-\Delta U)
=
\lambda^{\frac{p}{p-q}}(-\Delta U),
\]
while
\[
\lambda u^q
=
\lambda^{1+\frac{q}{p-q}}U^q
=
\lambda^{\frac{p}{p-q}}U^q
\]
and
\[
u^p
=
\lambda^{\frac{p}{p-q}}U^p.
\]
Thus the small-$\lambda$ regime is equivalent to the small-radius
regime
\[
R\to0^+.
\]

For $A>0$, let $U_A$ denote the maximal positive solution of the
radial initial-value problem
\begin{equation}
\label{eq:shooting}
U_A''
+\frac{N-1}{r}U_A'
+U_A^q+U_A^p
=0,
\qquad
U_A(0)=A,
\qquad
U_A'(0)=0.
\end{equation}
Define the endpoint of its interval of positivity by
\begin{equation}
\label{eq:firstzero}
R(A)
:=
\inf\{r>0:U_A(r)=0\},
\end{equation}
with the convention $R(A)=+\infty$ if $U_A$ remains positive for all
$r>0$.

For every $0<r<R(A)$, integration of~\eqref{eq:shooting} gives
\begin{equation}
\label{eq:monotone}
-r^{N-1}U_A'(r)
=
\int_0^r
s^{N-1}
\left(U_A(s)^q+U_A(s)^p\right)\,ds>0.
\end{equation}
Hence every positive shooting profile is strictly decreasing on
$(0,R(A))$. Consequently, positive radial solutions
of~\eqref{eq:scaled} are in one-to-one correspondence with values
$A>0$ satisfying
\[
R(A)=R.
\]

%%%%%%%%%%%%%%%%%%%%%%%%%%%%%%%%%%%%%%%%%%%%%%%%%%%%%%%%%%%%%%%%%
  
%%%%%%%%%%%%%%%%%%%%%%%%%%%%%%%%%%%%%%%%%%%%%%%%%%%%%%%%%%%%%%%%%
\section{The small-amplitude branch}
%%%%%%%%%%%%%%%%%%%%%%%%%%%%%%%%%%%%%%%%%%%%%%%%%%%%%%%%%%%%%%%%%



%%%%%%%%%%%%%

Define
\begin{equation}
\label{eq:Astar}
A_*
:=
\left(
\frac{1-q}{p-1}
\right)^{\frac1{p-q}}.
\end{equation}

\begin{lemma}
\label{lem:ratio}
The function
\[
s\longmapsto
\frac{s^q+s^p}{s}
=
s^{q-1}+s^{p-1}
\]
is strictly decreasing on $(0,A_*)$.
\end{lemma}

\begin{proof}
A direct computation gives
\[
\frac{d}{ds}
\left(
s^{q-1}+s^{p-1}
\right)
=
s^{q-2}
\left[
-(1-q)+(p-1)s^{p-q}
\right].
\]
The expression in brackets is negative exactly when
\[
s^{p-q}<\frac{1-q}{p-1},
\]
that is, for $0<s<A_*$.
\end{proof}

\begin{lemma}
\label{lem:picone-comparison}
Let $\Omega\subset\mathbb R^N$ be a bounded domain and let
$I=(0,M)$ for some $M>0$. Assume that
\[
f:I\longrightarrow(0,\infty)
\]
is continuous and that
\[
s\longmapsto\frac{f(s)}{s}
\]
is strictly decreasing on $I$.

Suppose that $u,v\in H_0^1(\Omega)\cap C(\overline\Omega)$ are
positive weak solutions of
\[
-\Delta w=f(w)
\qquad\text{in }\Omega,
\]
and that
\[
0<u(x),v(x)<M
\qquad\text{for every }x\in\Omega.
\]
Assume moreover that
\[
f(u)u,\ f(v)v\in L^1(\Omega).
\]
Then
\[
u\equiv v.
\]
\end{lemma}

\begin{proof}
For $\varepsilon>0$, Picone's identity applied to $u$ and
$v+\varepsilon$ gives
\[
|\nabla u|^2
-
\nabla(v+\varepsilon)\cdot
\nabla\left(\frac{u^2}{v+\varepsilon}\right)
=
\left|
\nabla u-
\frac{u}{v+\varepsilon}\nabla v
\right|^2
\geq0.
\]
Since
\[
\nabla(v+\varepsilon)=\nabla v,
\]
integration yields
\begin{equation}
\label{eq:picone-eps-1}
\int_\Omega |\nabla u|^2\,dx
\geq
\int_\Omega
\nabla v\cdot
\nabla\left(\frac{u^2}{v+\varepsilon}\right)\,dx.
\end{equation}

The function
\[
\frac{u^2}{v+\varepsilon}
\]
belongs to $H_0^1(\Omega)$ for every fixed $\varepsilon>0$, and
therefore it is an admissible test function in the weak equation for
$v$. Consequently,
\[
\int_\Omega
\nabla v\cdot
\nabla\left(\frac{u^2}{v+\varepsilon}\right)\,dx
=
\int_\Omega
f(v)\frac{u^2}{v+\varepsilon}\,dx.
\]
Testing the equation for $u$ with $u$ gives
\[
\int_\Omega|\nabla u|^2\,dx
=
\int_\Omega f(u)u\,dx.
\]
Thus
\begin{equation}
\label{eq:picone-eps-2}
\int_\Omega f(u)u\,dx
\geq
\int_\Omega
f(v)\frac{u^2}{v+\varepsilon}\,dx.
\end{equation}

Since the integrand on the right-hand side is nonnegative and
\[
f(v)\frac{u^2}{v+\varepsilon}
\longrightarrow
\frac{f(v)}{v}u^2
\qquad\text{a.e. in }\Omega
\]
as $\varepsilon\to0^+$, Fatou's lemma gives
\[
\int_\Omega
\frac{f(v)}{v}u^2\,dx
\leq
\liminf_{\varepsilon\to0^+}
\int_\Omega
f(v)\frac{u^2}{v+\varepsilon}\,dx
\leq
\int_\Omega f(u)u\,dx.
\]
Hence
\begin{equation}
\label{eq:picone-limit-1}
\int_\Omega
\left[
\frac{f(u)}{u}
-
\frac{f(v)}{v}
\right]u^2\,dx
\geq0.
\end{equation}

Interchanging $u$ and $v$, we likewise obtain
\begin{equation}
\label{eq:picone-limit-2}
\int_\Omega
\left[
\frac{f(v)}{v}
-
\frac{f(u)}{u}
\right]v^2\,dx
\geq0.
\end{equation}
Adding~\eqref{eq:picone-limit-1} and
\eqref{eq:picone-limit-2}, we obtain
\begin{equation}
\label{eq:picone-final}
\int_\Omega
\left[
\frac{f(u)}{u}
-
\frac{f(v)}{v}
\right]
(u^2-v^2)\,dx
\geq0.
\end{equation}

On the other hand, since $f(s)/s$ is strictly decreasing,
\[
\left[
\frac{f(u)}{u}
-
\frac{f(v)}{v}
\right](u-v)\leq0
\qquad\text{a.e. in }\Omega,
\]
with strict inequality wherever $u\neq v$. Since
\[
u^2-v^2=(u-v)(u+v)
\]
and $u+v>0$, it follows that
\[
\left[
\frac{f(u)}{u}
-
\frac{f(v)}{v}
\right]
(u^2-v^2)
\leq0
\qquad\text{a.e. in }\Omega,
\]
with strict inequality wherever $u\neq v$.

Together with~\eqref{eq:picone-final}, this forces
\[
u=v
\qquad\text{a.e. in }\Omega.
\]
By continuity,
\[
u\equiv v
\qquad\text{in }\overline\Omega.
\]
\end{proof}



\begin{proposition}
\label{prop:smalluniqueness}
For every $R>0$, problem~\eqref{eq:scaled} has at most one positive
solution $U$ satisfying
\[
\norm{U}_{L^\infty(B_R)}<A_*.
\]
\end{proposition}

\begin{proof}
Let $U_1,U_2$ be two such solutions and set
\[
f(s):=s^q+s^p.
\]
By Lemma~\ref{lem:ratio}, the map
\[
s\longmapsto \frac{f(s)}{s}
\]
is strictly decreasing on the range of both solutions.

\textcolor{red}{
Moreover,
\[
0<U_i(x)<A_*
\qquad
(x\in B_R,\ i=1,2),
\]
and
\[
f(U_i)U_i
=
U_i^{q+1}+U_i^{p+1}\in L^1(B_R),
\]
since $U_i$ is bounded. Therefore all the hypotheses of
Lemma~\ref{lem:picone-comparison} are satisfied with
\[
\Omega=B_R,
\qquad
M=A_*.
\]
It follows directly that
\[
U_1\equiv U_2.
\]
}
\end{proof}



%%%%%%%%%%%


We next identify the first-zero behavior as $A\to0^+$.

\begin{lemma}
\label{lem:rhoq}
\textcolor{red}{Assume that $0<q<1$. Let $V$ be the positive solution,
up to its first zero, of}
\begin{equation}
\label{eq:pureconcave}
V''
+\frac{N-1}{s}V'
+V^q
=0,
\qquad
V(0)=1,
\qquad
V'(0)=0.
\end{equation}
\textcolor{red}{If $\rho_q$ denotes its first positive zero, then}
\[
0<\rho_q\textcolor{red}{\leq
\sqrt{\frac{2N}{1-q}}}<\infty,
\qquad
V'(\rho_q)<0.
\]
\textcolor{red}{In particular, the first zero is transversal.}
\end{lemma}

\begin{proof}
As long as $V>0$, \textcolor{red}{equation~\eqref{eq:pureconcave}
can be written as}
\[
\textcolor{red}{\bigl(s^{N-1}V'(s)\bigr)'
=
-s^{N-1}V(s)^q.}
\]
\textcolor{red}{Using $V'(0)=0$, we obtain}
\[
-V'(s)
=
s^{1-N}
\int_0^s t^{N-1}V(t)^q\,dt>0.
\]
Thus $V$ is strictly decreasing on its interval of positivity.
Consequently, for
\[
\textcolor{red}{0\leq t\leq s},
\]
\[
V(t)\geq V(s),
\]
and hence
\[
\begin{aligned}
-V'(s)
&=
\textcolor{red}{s^{1-N}
\int_0^s t^{N-1}V(t)^q\,dt}\\
&\geq
s^{1-N}V(s)^q
\int_0^s t^{N-1}\,dt
=
\frac{s}{N}V(s)^q.
\end{aligned}
\]
Therefore
\[
-\frac{V'(s)}{V(s)^q}
\geq
\frac{s}{N}.
\]
\textcolor{red}{Since $0<q<1$,
\[
\frac{d}{ds}V(s)^{1-q}
=
(1-q)V(s)^{-q}V'(s),
\]
and hence
\[
\frac{d}{ds}V(s)^{1-q}
\leq
-\frac{1-q}{N}s.
\]}
Integrating from $0$ to $s$ \textcolor{red}{and using $V(0)=1$} gives
\[
V(s)^{1-q}
\leq
1-\frac{1-q}{2N}s^2.
\]
\textcolor{red}{If $V$ remained positive up to
\[
s_*:=\sqrt{\frac{2N}{1-q}},
\]
then the preceding inequality would give
\[
0<V(s_*)^{1-q}\leq0,
\]
which is impossible. Therefore $V$ has a first positive zero
$\rho_q$ satisfying
\[
0<\rho_q\leq
\sqrt{\frac{2N}{1-q}}<\infty.
\]}

Finally, \textcolor{red}{passing to the limit $s\uparrow\rho_q$ in the
integral identity} gives
\[
-\rho_q^{N-1}V'(\rho_q)
=
\int_0^{\rho_q}
t^{N-1}V(t)^q\,dt.
\]
\textcolor{red}{The integral on the right-hand side is strictly
positive, since $V(t)>0$ for $0\leq t<\rho_q$.} Hence
\[
V'(\rho_q)<0.
\]
\textcolor{red}{Thus the first zero is transversal.}
\end{proof}


\begin{proposition}
\label{prop:smallasymptotic}
\textcolor{red}{Assume that $0<q<1<p$. Then,} as $A\to0^+$,
\begin{equation}
\label{eq:Rsmall}
R(A)
=
\rho_q
A^{\frac{1-q}{2}}
\bigl(1+o(1)\bigr).
\end{equation}
\end{proposition}

\begin{proof}
Set
\[
U_A(r)
=
A V_A(s),
\qquad
s=A^{-\frac{1-q}{2}}r.
\]
Then $V_A$ solves
\begin{equation}
\label{eq:VA-small}
V_A''
+\frac{N-1}{s}V_A'
+V_A^q
+A^{p-q}V_A^p
=0,
\qquad
V_A(0)=1,
\qquad
V_A'(0)=0.
\end{equation}
Let $\rho_A$ denote the first positive zero of $V_A$.

\textcolor{red}{We first prove
\[
\liminf_{A\to0^+}\rho_A\geq\rho_q.
\]}
Fix \textcolor{red}{$\delta\in(0,\rho_q)$ and set
\[
m_\delta
:=
\min_{0\leq s\leq\rho_q-\delta}V(s)>0.
\]}
On a \textcolor{red}{compact} neighborhood of
\textcolor{red}{$[m_\delta,1]$ contained in $(0,\infty)$}, the
\textcolor{red}{maps $t\mapsto t^q$ and $t\mapsto t^p$} are locally
Lipschitz. \textcolor{red}{Since $A^{p-q}\to0$, continuous dependence
for the radial initial-value problem, equivalently from its integral
formulation, yields}
\[
V_A\longrightarrow V
\qquad\text{in }C^1([0,\rho_q-\delta])
\]
as $A\to0^+$. \textcolor{red}{Hence, for all sufficiently small $A$,
\[
V_A(s)\geq\frac{m_\delta}{2}>0
\qquad
(0\leq s\leq\rho_q-\delta),
\]
and therefore
\[
\rho_A>\rho_q-\delta.
\]
Since $\delta>0$ is arbitrary,}
\[
\liminf_{A\to0^+}\rho_A\geq\rho_q.
\]

We next prove the reverse inequality. \textcolor{red}{Since}
\[
V(\rho_q)=0
\qquad\text{and}\qquad
V'(\rho_q)<0,
\]
\textcolor{red}{choose $c_0>0$ and $\eta>0$ such that
\[
V'(s)\leq-2c_0
\qquad
\text{for }
s\in[\rho_q-\eta,\rho_q].
\]}

Fix $\varepsilon>0$.
\textcolor{red}{Choose}
\[
s_0\in
\textcolor{red}{
\left(
\rho_q-\min\{\eta,\varepsilon/2\},
\rho_q
\right)}
\]
so close to $\rho_q$ that
\[
\textcolor{red}{
V(s_0)
<
\frac{c_0\varepsilon}{4}
\left(
\frac{s_0}{s_0+\varepsilon/2}
\right)^{N-1}.}
\]
By continuous dependence
\textcolor{red}{on the fixed interval $[0,s_0]$},
\[
V_A(s_0)\to V(s_0),
\qquad
V_A'(s_0)\to V'(s_0).
\]
\textcolor{red}{Thus, for all sufficiently small $A$,}
\[
V_A'(s_0)\leq -c_0
\]
\textcolor{red}{and
\[
V_A(s_0)
<
\frac{c_0\varepsilon}{2}
\left(
\frac{s_0}{s_0+\varepsilon/2}
\right)^{N-1}.
\]}

As long as $V_A$ remains positive, equation~\eqref{eq:VA-small}
implies
\[
\bigl(s^{N-1}V_A'(s)\bigr)'
=
-s^{N-1}
\left(
V_A(s)^q+A^{p-q}V_A(s)^p
\right)
<0.
\]
Therefore, for $s\geq s_0$, as long as $V_A(s)>0$,
\[
s^{N-1}V_A'(s)
\leq
s_0^{N-1}V_A'(s_0),
\]
and hence
\[
V_A'(s)
\leq
-c_0\left(\frac{s_0}{s}\right)^{N-1}.
\]
\textcolor{red}{Suppose, by contradiction, that $V_A$ remains positive
on $[s_0,s_0+\varepsilon/2]$. Then
\[
V_A'(s)
\leq
-c_0
\left(
\frac{s_0}{s_0+\varepsilon/2}
\right)^{N-1}
\]
throughout that interval. Hence
\[
\begin{aligned}
V_A(s_0+\varepsilon/2)
&\leq
V_A(s_0)
-\frac{c_0\varepsilon}{2}
\left(
\frac{s_0}{s_0+\varepsilon/2}
\right)^{N-1}\\
&<0,
\end{aligned}
\]
a contradiction. Therefore
\[
\rho_A<s_0+\frac{\varepsilon}{2}<\rho_q+\varepsilon.
\]}
It follows that
\[
\limsup_{A\to0^+}\rho_A\leq\rho_q.
\]
Combining the two inequalities, we conclude that
\[
\rho_A\longrightarrow\rho_q.
\]

Returning to the original radial variable,
\[
R(A)
=
A^{\frac{1-q}{2}}\rho_A,
\]
and therefore
\[
R(A)
=
\rho_q A^{\frac{1-q}{2}}
\bigl(1+o(1)\bigr),
\]
which proves~\eqref{eq:Rsmall}.
\end{proof}


\begin{corollary}
\label{cor:smallbranch}
There exist $R_0>0$ and $A_0\in(0,A_*)$ such that, for every
$R\in(0,R_0)$, problem~\eqref{eq:scaled} admits a positive radial
solution \textcolor{red}{$U$ satisfying}
\[
\textcolor{red}{U(0)<A_0.}
\]
\textcolor{red}{Moreover,
\[
\|U\|_{L^\infty(B_R)}=U(0)<A_0<A_*.
\]}
\textcolor{red}{Proposition~\ref{prop:smalluniqueness} applies to
all positive solutions of~\eqref{eq:scaled}, then} this solution is
unique among all positive solutions satisfying
\[
\norm{U}_{L^\infty(B_R)}<A_*.
\]
\end{corollary}

\begin{proof}
By Proposition~\ref{prop:smallasymptotic},
\textcolor{red}{
\[
R(A)
=
\rho_q A^{\frac{1-q}{2}}(1+o(1))
\qquad\text{as }A\to0^+.
\]
In particular,}
\[
R(A)\to0
\qquad\text{as }A\to0^+.
\]

The transversality \textcolor{red}{of the first zero}, together with
continuous dependence of the radial initial-value problem
\textcolor{red}{on the initial value, implies persistence and
continuous dependence of the first zero. Hence there exists
$\delta>0$ such that}
\[
A\longmapsto R(A)
\]
is continuous \textcolor{red}{on $(0,\delta)$}.

\textcolor{red}{Choose
\[
A_0\in(0,\min\{\delta,A_*\})
\]
and set
\[
R_0:=R(A_0)>0.
\]
Fix $R\in(0,R_0)$. Since
\[
R(A)\to0
\qquad\text{as }A\to0^+,
\]
there exists $A_1\in(0,A_0)$ such that
\[
R(A_1)<R.
\]
On the other hand,
\[
R(A_0)=R_0>R.
\]
By continuity of the first-zero map, there exists
\[
A\in(A_1,A_0)
\]
such that
\[
R(A)=R.
\]}
The corresponding shooting profile $U_A$ is therefore a positive
radial solution of~\eqref{eq:scaled}
\textcolor{red}{in $B_R$}.

\textcolor{red}{Moreover,} every positive radial shooting profile is
strictly decreasing \textcolor{red}{on its interval of positivity}.
Hence
\[
\norm{U_A}_{L^\infty(B_R)}
=
U_A(0)
\textcolor{red}{=A<A_0}<A_*.
\]

Finally, \textcolor{red}{if Proposition~\ref{prop:smalluniqueness}
gives uniqueness among all positive solutions satisfying
\[
\|U\|_{L^\infty(B_R)}<A_*,
\]
then the asserted uniqueness follows directly.}
\end{proof}


%%%%%%%%%%%%%%%%%%%%%%%%%%%%%%%%%%%%%%%%%%%%%%%%%%%%%%%%%%%%%%%%%
  
%%%%%%%%%%%%%%%%%%%%%%%%%%%%%%%%%%%%%%%%%%%%%%%%%%%%%%%%%%%%%%%%%
\section{The supercritical Lane--Emden orbit}
%%%%%%%%%%%%%%%%%%%%%%%%%%%%%%%%%%%%%%%%%%%%%%%%%%%%%%%%%%%%%%%%%

We now study the radial initial-value problem
\begin{equation}
\label{eq:LE}
v_0''
+\frac{N-1}{r}v_0'
+v_0^p
=0,
\qquad
v_0(0)=1,
\qquad
v_0'(0)=0.
\end{equation}

Set
\begin{equation}
\label{eq:aAB}
a:=\frac{2}{p-1},
\qquad
A:=N-2-2a,
\qquad
B:=a(N-2-a).
\end{equation}
Since $p>\pc$, we have
\[
0<a<\frac{N-2}{2},
\]
and therefore
\[
A>0,
\qquad
B>0.
\]

Introduce the Emden--Fowler variables
\[
t=\log r,
\qquad
x_0(t)=r^a v_0(r).
\]
Then $x_0$ satisfies
\begin{equation}
\label{eq:EF0}
x_0''
+Ax_0'
-Bx_0
+x_0^p
=0.
\end{equation}
The unique positive equilibrium of~\eqref{eq:EF0} is
\begin{equation}
\label{eq:ell}
\ell
:=
B^{\frac1{p-1}}.
\end{equation}

\begin{lemma}
\label{lem:LEasymptotics}
\textcolor{red}{Assume that $N\geq3$ and
\[
p>\pc:=\frac{N+2}{N-2}.
\]}
The solution $v_0$ of~\eqref{eq:LE} is positive on $(0,\infty)$ and
\[
x_0(t)\longrightarrow\ell,
\qquad
x_0'(t)\longrightarrow0
\qquad
\text{as }t\to+\infty.
\]
Equivalently,
\begin{equation}
\label{eq:LE-tail}
r^{\frac2{p-1}}v_0(r)
\longrightarrow
\left[
\frac2{p-1}
\left(
N-2-\frac2{p-1}
\right)
\right]^{\frac1{p-1}}
\qquad
\text{as }r\to+\infty.
\end{equation}
\end{lemma}

\begin{proof}
Suppose that $v_0$ had a first zero at some finite radius $R>0$.
Since $v_0>0$ in $(0,R)$, it would define a nontrivial positive
solution of
\[
-\Delta v=v^p
\quad\text{in }B_R,
\qquad
v=0
\quad\text{on }\partial B_R.
\]
Since \textcolor{red}{$p>\pc$} and $B_R$ is star-shaped, the Pohozaev
identity excludes the existence of such a solution. Hence
\[
v_0(r)>0
\qquad
\text{for every }r>0.
\]

For~\eqref{eq:EF0}, consider the Lyapunov function
\[
E(x,y)
:=
\frac12y^2
-\frac{B}{2}x^2
+\frac1{p+1}x^{p+1}.
\]
Along the orbit $(x_0,x_0')$,
\[
\frac{d}{dt}E(x_0(t),x_0'(t))
=
-A(x_0'(t))^2
\leq0.
\]

\textcolor{red}{Regularity of the radial solution at the origin gives
\[
v_0(r)=1+O(r^2),
\qquad
v_0'(r)=O(r)
\qquad
\text{as }r\to0^+.
\]
Consequently,}
\[
x_0(t)=e^{at}\bigl(1+o(1)\bigr),
\qquad
x_0'(t)=ae^{at}\bigl(1+o(1)\bigr)
\qquad
\text{as }t\to-\infty.
\]
Consequently,
\[
E(x_0(t),x_0'(t))\longrightarrow0
\qquad
\text{as }t\to-\infty.
\]
Moreover,
\[
E(x_0(t),x_0'(t))
=
-A\int_{-\infty}^{t}(x_0'(s))^2\,ds.
\]
\textcolor{red}{Since
\[
x_0'(s)\sim ae^{as}>0
\qquad\text{as }s\to-\infty,
\]
the integral on the right-hand side is strictly positive for every
finite $t$.}
It follows that
\[
E(x_0(t),x_0'(t))<0
\qquad
\text{for every finite }t.
\]

We next show that the positive semi-orbit is bounded. Since
\[
E(x_0,x_0')
=
\frac12(x_0')^2
-\frac{B}{2}x_0^2
+\frac1{p+1}x_0^{p+1}
<0,
\]
we obtain
\[
\frac1{p+1}x_0^{p+1}
<
\frac{B}{2}x_0^2.
\]
\textcolor{red}{Since $x_0>0$, this yields the explicit bound
\[
0<x_0(t)
<
\left(\frac{B(p+1)}{2}\right)^{\frac1{p-1}}
\qquad\text{for every }t\in\mathbb R.
\]}
Thus $x_0$ is uniformly bounded.

\textcolor{red}{Moreover, again using $E<0$,
\[
\frac12(x_0')^2
<
\frac{B}{2}x_0^2
-\frac1{p+1}x_0^{p+1}
\leq
\frac{B}{2}x_0^2,
\]
and therefore
\[
|x_0'(t)|<\sqrt{B}\,x_0(t).
\]}
Hence $x_0'$ is uniformly bounded as well, and the positive
semi-orbit is \textcolor{red}{relatively compact in the phase plane}.

Since $E$ is nonincreasing and bounded below along the orbit, there
exists
\[
E_\infty
:=
\lim_{t\to+\infty}E(x_0(t),x_0'(t)).
\]
For every fixed $t_0\in\R$,
\[
E_\infty\leq E(x_0(t_0),x_0'(t_0))<0,
\]
so
\[
E_\infty<0.
\]

\textcolor{red}{To apply LaSalle's invariance principle without any
difficulty at $x=0$ when $p$ is not an integer, extend the vector
field to the whole phase plane by replacing $x^p$ with
$|x|^{p-1}x$. Thus consider
\[
x'=y,
\qquad
y'=-Ay+Bx-|x|^{p-1}x.
\]
This system agrees with~\eqref{eq:EF0} along the positive orbit.
The corresponding globally defined Lyapunov function is
\[
\widetilde E(x,y)
=
\frac12y^2
-\frac{B}{2}x^2
+\frac1{p+1}|x|^{p+1},
\]
and along its trajectories
\[
\frac{d}{dt}\widetilde E(x(t),y(t))
=
-Ay(t)^2.
\]}

Because the positive semi-orbit is \textcolor{red}{relatively compact},
its $\omega$-limit set is nonempty, compact, connected, and invariant.
By LaSalle's invariance principle it is contained in the largest
invariant subset of
\[
\{(x,y):y=0\}.
\]
\textcolor{red}{Since the $\omega$-limit set is contained in the
nonnegative half-plane,} invariance then forces
\[
Bx-x^p=0.
\]
Thus the only nonnegative invariant points in this set are the
equilibria
\[
(0,0)
\qquad\text{and}\qquad
(\ell,0).
\]

The equilibrium $(0,0)$ has energy $0$, whereas the energy is constant
and equal to $E_\infty<0$ on the $\omega$-limit set. Hence $(0,0)$
cannot belong to it. Since the $\omega$-limit set is nonempty, we
conclude that
\[
\omega(x_0,x_0')=\{(\ell,0)\}.
\]
Consequently,
\[
x_0(t)\longrightarrow\ell,
\qquad
x_0'(t)\longrightarrow0
\qquad
\text{as }t\to+\infty.
\]

\textcolor{red}{Finally, since
\[
\ell=B^{\frac1{p-1}}
=
\left[
\frac2{p-1}
\left(
N-2-\frac2{p-1}
\right)
\right]^{\frac1{p-1}},
\]
the preceding convergence is equivalent to~\eqref{eq:LE-tail}.}
This proves the assertion.
\end{proof}

The equilibrium $(\ell,0)$ is exponentially asymptotically stable.
Indeed, writing $x=\ell+z$ and linearizing~\eqref{eq:EF0} at
$(\ell,0)$ gives
\[
z''
+Az'
+(p-1)Bz
=0.
\]
Its characteristic polynomial is
\[
\mu^2+A\mu+(p-1)B.
\]
Since
\[
A>0
\qquad\text{and}\qquad
(p-1)B>0,
\]
both characteristic roots have strictly negative real parts.
\textcolor{red}{Hence the corresponding first-order linearized matrix
is Hurwitz, and $(\ell,0)$ is a hyperbolic exponentially
asymptotically stable equilibrium of the autonomous
Emden--Fowler system.}

We record the following robust stability property, which will be
used in the large-amplitude analysis.

\begin{lemma}
\label{lem:robust}
There exist numbers $\eta>0$, $\eta_0>0$, and $\delta>0$ with
\[
0<\eta<\frac{\ell}{4}
\]
such that the following holds.

Let $d:[T,T_1]\to[0,\delta]$ be continuous and let $(x,y)$ be a
solution, as long as $x>0$, of
\begin{equation}
\label{eq:perturbed-system}
\begin{cases}
x'=y,\\
y'=-Ay+Bx-x^p-d(t)x^q,
\end{cases}
\end{equation}
with
\[
|x(T)-\ell|+|y(T)|<\eta_0.
\]
Then
\[
|x(t)-\ell|+|y(t)|<\eta
\qquad
\text{for every }t\in[T,T_1].
\]
In particular,
\[
x(t)>\frac{\ell}{2}
\qquad
(T\leq t\leq T_1).
\]
\end{lemma}

\begin{proof}
Set
\[
Z=
\binom{x-\ell}{y}.
\]
\textcolor{red}{In what follows we use the $\ell^1$-norm
\[
|Z|:=|x-\ell|+|y|,
\]
together with its induced matrix norm.}
Since $\ell>0$, \textcolor{red}{we shall choose} $\eta<\ell/4$.
Thus, whenever
$|Z|\leq\eta$, one has
\[
x>\ell-\eta>\frac{3\ell}{4}>0,
\]
so the powers $x^p$ and $x^q$ are smooth functions of $x$ in the
neighborhood under consideration.

Near $(\ell,0)$, system~\eqref{eq:perturbed-system} can therefore be
written as
\[
Z'
=
MZ+\mathcal N(Z)+d(t)\mathcal F(Z),
\]
where
\[
M=
\begin{pmatrix}
0&1\\
-(p-1)B&-A
\end{pmatrix},
\]
\textcolor{red}{and, writing $Z=(z,y)^T$ with $z=x-\ell$, one may take
\[
\mathcal N(Z)
=
\binom{0}{
B(\ell+z)-(\ell+z)^p+(p-1)Bz
},
\qquad
\mathcal F(Z)
=
\binom{0}{-(\ell+z)^q}.
\]
Indeed, since $\ell^{p-1}=B$,
\[
B\ell-\ell^p=0,
\]
and
\[
\left.
\frac{d}{dz}
\left[
B(\ell+z)-(\ell+z)^p
\right]
\right|_{z=0}
=
-(p-1)B.
\]
Therefore Taylor's theorem gives, after possibly decreasing $\eta$,}
\[
|\mathcal N(Z)|
\leq C|Z|^2,
\qquad
|\mathcal F(Z)|
\leq C
\]
\textcolor{red}{for every $|Z|\leq\eta$, with a constant $C>0$
independent of $d$, $T$, and $T_1$.}

Since $M$ is Hurwitz, there exist constants $K,\omega>0$ such that
\[
\norm{e^{Mt}}
\leq
Ke^{-\omega t}
\qquad
(t\geq0).
\]
\textcolor{red}{Here $K$ and $\omega$ depend only on the fixed matrix
$M$, hence only on $N$ and $p$.}
By the variation-of-constants formula, as long as $|Z(s)|\leq\eta$
for $s\in[T,t]$,
\begin{align*}
|Z(t)|
&\leq
Ke^{-\omega(t-T)}|Z(T)|
\\
&\quad
+
K\int_T^t
e^{-\omega(t-s)}
\left(
C|Z(s)|^2
+C\delta
\right)\,ds
\\
&\leq
K|Z(T)|
+
\frac{KC\eta}{\omega}
\sup_{T\leq s\leq t}|Z(s)|
+
\frac{KC\delta}{\omega}.
\end{align*}

Choose $\eta>0$ sufficiently small that
\[
\textcolor{red}{
0<\eta<\frac{\ell}{4}
\qquad\text{and}\qquad
}
\frac{KC\eta}{\omega}<\frac14,
\]
and then choose $\eta_0>0$ and $\delta>0$ sufficiently small that
\[
K\eta_0+\frac{KC\delta}{\omega}
<
\frac{\eta}{2}.
\]
\textcolor{red}{We may also assume $\eta_0<\eta$, so that the initial
point already lies strictly inside the neighborhood
$\{|Z|<\eta\}$.}

Suppose, by contradiction, that there exists a first exit time
$t_*\in(T,T_1]$ such that
\[
|Z(t_*)|=\eta
\]
and
\[
|Z(t)|<\eta
\qquad
(T\leq t<t_*).
\]
\textcolor{red}{By continuity, $|Z(t)|\leq\eta$ on $[T,t_*]$, so the
preceding variation-of-constants estimate is valid throughout this
interval.}
Applying the previous estimate on $[T,t_*]$ gives
\[
\sup_{T\leq t\leq t_*}|Z(t)|
<
\frac{\eta}{2}
+
\frac14
\sup_{T\leq t\leq t_*}|Z(t)|.
\]
Hence
\[
\sup_{T\leq t\leq t_*}|Z(t)|
<
\frac{2\eta}{3}
<\eta,
\]
contradicting the definition of $t_*$. Therefore
\[
|Z(t)|<\eta
\qquad
\text{for every }t\in[T,T_1].
\]
In particular,
\[
x(t)>\ell-\eta>\frac{3\ell}{4}>\frac{\ell}{2},
\]
which completes the proof.
\end{proof}
    


%%%%%%%%%%%%%%%%%%%%%%%%%%%%%%%%%%%%%%%%%%%%%%%%%%%%%%%%%%%%%%%%%
%%%%%%%%%%%%%%%%%%%%%%%%%%%%%%%%%%%%%%%%%%%%%%%%%%%%%%%%%%%%%%%%%
\section{Large-amplitude exclusion}
%%%%%%%%%%%%%%%%%%%%%%%%%%%%%%%%%%%%%%%%%%%%%%%%%%%%%%%%%%%%%%%%%

For $\varepsilon>0$, let $v_\varepsilon$ denote the maximal positive
solution of
\begin{equation}
\label{eq:veps}
v_\varepsilon''
+\frac{N-1}{r}v_\varepsilon'
+v_\varepsilon^p
+\varepsilon v_\varepsilon^q
=0,
\qquad
v_\varepsilon(0)=1,
\qquad
v_\varepsilon'(0)=0.
\end{equation}
Let $S(\varepsilon)$ denote its first positive zero, with the
convention
\(
S(\varepsilon)=+\infty
\)
if $v_\varepsilon$ remains positive for all $r>0$.

\begin{lemma}
\label{lem:S-lower}
\textcolor{red}{Assume that $0<q<1<p$ and $p>\pc$.}
There exist constants $c>0$ and $\varepsilon_0>0$ such that
\begin{equation}
\label{eq:Sbound}
S(\varepsilon)
\geq
c\,
\varepsilon^{-\frac{p-1}{2(p-q)}}
\end{equation}
for every $0<\varepsilon<\varepsilon_0$.
\end{lemma}

\begin{proof}
Use the Emden--Fowler variables
\[
t=\log r,
\qquad
x_\varepsilon(t)=r^av_\varepsilon(r),
\qquad
a=\frac2{p-1}.
\]
As long as $v_\varepsilon$ remains positive, a direct computation
gives
\begin{equation}
\label{eq:EFeps}
x_\varepsilon''
+Ax_\varepsilon'
-Bx_\varepsilon
+x_\varepsilon^p
+\varepsilon e^{mt}x_\varepsilon^q
=0,
\end{equation}
where
\begin{equation}
\label{eq:m}
m
=
2+a(1-q)
=
\frac{2(p-q)}{p-1}
\textcolor{red}{>0}.
\end{equation}

By Lemma~\ref{lem:LEasymptotics}, choose $T$ so large that
\[
|x_0(T)-\ell|+|x_0'(T)|
<
\frac{\eta_0}{2},
\]
where $\eta_0$ is supplied by Lemma~\ref{lem:robust}.

Since $v_0$ is positive on the compact interval $[0,e^T]$, set
\[
\kappa_T:=\min_{0\le r\le e^T}v_0(r)>0.
\]
On the region $v\ge \kappa_T/2$, the map
\[
(v,\varepsilon)\longmapsto v^p+\varepsilon v^q
\]
is locally Lipschitz in $v$, uniformly for $\varepsilon$ in a bounded
interval. The radial equation may equivalently be written in integral
form as
\[
v(r)=1-\int_0^r t^{1-N}\int_0^t
s^{N-1}\bigl(v(s)^p+\varepsilon v(s)^q\bigr)\,ds\,dt,
\]
which removes the apparent singularity at $r=0$.
\textcolor{red}{Since $v_0\geq\kappa_T>0$ on $[0,e^T]$, continuous
dependence for this integral equation, together with a standard
continuation argument inside the region
$v\geq\kappa_T/2$, gives}
\[
v_\varepsilon\ge \frac{\kappa_T}{2}
\quad\hbox{on }[0,e^T],
\qquad
v_\varepsilon\longrightarrow v_0
\quad\hbox{in }C^1([0,e^T])
\]
\textcolor{red}{as $\varepsilon\to0^+$.}
In particular,
\[
(x_\varepsilon(T),x_\varepsilon'(T))
\longrightarrow
(x_0(T),x_0'(T))
\qquad
\text{as }\varepsilon\to0^+.
\]
Hence, after reducing $\varepsilon_0$ if necessary,
\[
|x_\varepsilon(T)-\ell|
+
|x_\varepsilon'(T)|
<
\eta_0.
\]

Choose $\tau_0<0$ so that
\[
\textcolor{red}{
e^{m\tau_0}<\frac{\delta}{2},
}
\]
where $\delta$ is given by Lemma~\ref{lem:robust}, and define
\[
T_\varepsilon
:=
-\frac1m\log\varepsilon+\tau_0.
\]
For $\varepsilon>0$ sufficiently small,
\[
T_\varepsilon>T.
\]
Moreover, for every $t\in[T,T_\varepsilon]$,
\[
0\leq
\varepsilon e^{mt}
\leq
\varepsilon e^{mT_\varepsilon}
=
e^{m\tau_0}
<
\textcolor{red}{\frac{\delta}{2}<\delta}.
\]

\textcolor{red}{If $S(\varepsilon)=+\infty$, there is nothing to
prove. We may therefore suppose that $S(\varepsilon)<+\infty$ and
set
\[
t_\varepsilon^*:=\log S(\varepsilon).
\]}
We claim that $x_\varepsilon$ remains positive throughout
$[T,T_\varepsilon]$.
\textcolor{red}{Suppose, by contradiction, that
\[
t_\varepsilon^*\leq T_\varepsilon.
\]
For every
\[
T_1\in[T,t_\varepsilon^*),
\]
the solution is positive on $[T,T_1]$, and hence
Lemma~\ref{lem:robust} applies on this interval with}
\[
d(t):=\varepsilon e^{mt}\in[0,\delta].
\]
The lemma gives
\[
x_\varepsilon(t)>\frac{\ell}{2}
\]
\textcolor{red}{for every $t\in[T,T_1]$. Since $T_1<t_\varepsilon^*$
is arbitrary,}
\[
x_\varepsilon(t)\geq\frac{\ell}{2}
\qquad
\textcolor{red}{\text{for }T\leq t<t_\varepsilon^*.}
\]
\textcolor{red}{Passing to the limit as
$t\uparrow t_\varepsilon^*$ yields
\[
x_\varepsilon(t_\varepsilon^*)
\geq\frac{\ell}{2}>0.
\]
On the other hand, by the definition of the first zero,
\[
x_\varepsilon(t_\varepsilon^*)
=
S(\varepsilon)^a
v_\varepsilon(S(\varepsilon))
=0,
\]
which is a contradiction. Therefore
\[
t_\varepsilon^*>T_\varepsilon,
\qquad\text{equivalently}\qquad
S(\varepsilon)>e^{T_\varepsilon}.
\]}

\textcolor{red}{Thus, in all cases,}
\[
S(\varepsilon)>e^{T_\varepsilon}.
\]
Finally,
\[
e^{T_\varepsilon}
=
e^{\tau_0}
\varepsilon^{-1/m}
=
e^{\tau_0}
\varepsilon^{-\frac{p-1}{2(p-q)}}.
\]
Thus~\eqref{eq:Sbound} holds with
\[
c:=e^{\tau_0}>0.
\]
\end{proof}

We now transfer this estimate to the original shooting family.

\begin{proposition}
\label{prop:largeA}
\textcolor{red}{Assume that $0<q<1<p$ and $p>\pc$.}
There exist $A_1>0$ and $c_*>0$ such that
\(
R(A)\geq c_*
\)
for every \(A\geq A_1\).
\end{proposition}

\begin{proof}
For $A>0$, set
\[
U_A(r)
=
A v_\varepsilon(s),
\qquad
s=A^{\frac{p-1}{2}}r,
\qquad
\varepsilon=A^{q-p}.
\]
Then $v_\varepsilon$ solves~\eqref{eq:veps}, and the corresponding
first zeros satisfy
\begin{equation}
\label{eq:R-S}
R(A)
=
A^{-\frac{p-1}{2}}
S(A^{q-p}).
\end{equation}
\textcolor{red}{Here the identity is understood in the extended
sense: if $S(A^{q-p})=+\infty$, then $R(A)=+\infty$.}

Since $q-p<0$, we have
\[
A^{q-p}\longrightarrow0
\qquad
\text{as }A\to+\infty.
\]
Choose $A_1>0$ so large that
\[
A^{q-p}<\varepsilon_0
\qquad
\text{for every }A\geq A_1.
\]
\textcolor{red}{Fix $A\geq A_1$. If
$S(A^{q-p})=+\infty$, then $R(A)=+\infty$ and the desired estimate
is immediate. We may therefore assume that
$S(A^{q-p})<+\infty$.}
By Lemma~\ref{lem:S-lower},
\begin{align*}
R(A)
&\geq
c
A^{-\frac{p-1}{2}}
\left(A^{q-p}\right)^{-\frac{p-1}{2(p-q)}}
\\
&=
c
A^{-\frac{p-1}{2}}
A^{\frac{p-1}{2}}
=
c.
\end{align*}
\textcolor{red}{Indeed, since $q-p=-(p-q)$,
\[
\left(A^{q-p}\right)^{-\frac{p-1}{2(p-q)}}
=
A^{\frac{p-1}{2}}.
\]}
Thus one may take
\[
c_*:=c.
\]
\end{proof}

Intermediate amplitudes are easier.

\begin{lemma}
\label{lem:intermediate}
For every $0<A_0<A_1<\infty$, there exists $c_{0,1}>0$ such that
\(
R(A)\geq c_{0,1}
\)
for every \(A\in[A_0,A_1]\).
\end{lemma}

\begin{proof}
Fix $A\in[A_0,A_1]$.
\textcolor{red}{If $R(A)=+\infty$, the desired estimate is immediate.
Thus suppose that $R(A)<+\infty$.}
As long as $U_A$ remains positive, it is
strictly decreasing, and therefore
\[
0<U_A(r)\leq A\leq A_1.
\]
Using~\eqref{eq:monotone}, for every $0<r<R(A)$ we obtain
\[
-U_A'(r)
=
r^{1-N}
\int_0^r
s^{N-1}
\left(
U_A(s)^q+U_A(s)^p
\right)\,ds
\leq
\frac{r}{N}
\left(A_1^q+A_1^p\right).
\]
Integrating from $0$ to $r<R(A)$ gives
\[
U_A(r)
\geq
A
-
\frac{r^2}{2N}
\left(A_1^q+A_1^p\right)
\geq
A_0
-
\frac{r^2}{2N}
\left(A_1^q+A_1^p\right).
\]

\textcolor{red}{For instance, choose}
\[
\textcolor{red}{
c_{0,1}
:=
\left(
\frac{N A_0}{A_1^q+A_1^p}
\right)^{1/2}>0.
}
\]
\textcolor{red}{Then}
\[
\frac{c_{0,1}^2}{2N}
\left(A_1^q+A_1^p\right)
=
\frac{A_0}{2}.
\]
We claim that
\[
R(A)\geq c_{0,1}.
\]
Indeed, if $R(A)<c_{0,1}$, then the previous estimate holds for every
$0<r<R(A)$. Letting $r\to R(A)^-$ and using continuity of $U_A$, we
would obtain
\[
0=U_A(R(A))
\geq
A_0
-
\frac{R(A)^2}{2N}
\left(A_1^q+A_1^p\right)
>
\frac{A_0}{2},
\]
a contradiction. Therefore
\[
R(A)\geq c_{0,1},
\]
uniformly for $A\in[A_0,A_1]$.
\end{proof}

%%%%%%%%%%%%%%%%%%%%%%%%%%%%%%%%%%%%%%%%%%%%%%%%%%%%%%%%%%%%%%%%%
  
%%%%%%%%%%%%%%%%%%%%%%%%%%%%%%%%%%%%%%%%%%%%%%%%%%%%%%%%%%%%%%%%%
\section{Radial uniqueness for small balls}
%%%%%%%%%%%%%%%%%%%%%%%%%%%%%%%%%%%%%%%%%%%%%%%%%%%%%%%%%%%%%%%%%

\begin{theorem}
\label{thm:radial}
There exists $R_*>0$ such that, for every
\(
0<R<R_*,
\)
problem~\eqref{eq:scaled} admits exactly one positive radial solution,
denoted by $U_R$. Moreover,
\[
\norm{U_R}_{L^\infty(B_R)}<A_*.
\]
\end{theorem}

\begin{proof}
\textcolor{red}{By Corollary~\ref{cor:smallbranch}, choose
\[
0<A_0<A_*
\]
and $R_0>0$ such that, for every $0<R<R_0$,
problem~\eqref{eq:scaled} admits a positive radial solution whose
center value satisfies}
\[
U(0)<A_0.
\]
\textcolor{red}{In particular, by
Proposition~\ref{prop:smallasymptotic}, the corresponding shooting
parameter tends to zero along the small-amplitude branch as
$R\to0^+$.}

Let $A_1>A_0$ be large enough so that
Proposition~\ref{prop:largeA} applies. By
Lemma~\ref{lem:intermediate}, there exists $c_{0,1}>0$ such that
\[
R(A)\geq c_{0,1}
\qquad
\text{for every }A\in[A_0,A_1],
\]
whereas Proposition~\ref{prop:largeA} yields
\[
R(A)\geq c_*
\qquad
\text{for every }A\geq A_1.
\]
Consequently,
\[
R(A)\geq
\min\{c_{0,1},c_*\}>0
\qquad
\text{for every }A\geq A_0.
\]

Choose $R_*>0$ such that
\[
\textcolor{red}{
R_*<
\min\{R_0,c_{0,1},c_*\}.
}
\]
\textcolor{red}{Then the small-amplitude existence conclusion holds
for every $0<R<R_*$, while at the same time}
\[
R_*<\min\{c_{0,1},c_*\}.
\]

Let now $0<R<R_*$ and let $U$ be any positive radial solution
of~\eqref{eq:scaled}. By the shooting formulation,
\[
U(0)=A
\]
for some $A>0$ satisfying
\[
R(A)=R.
\]
If $A\geq A_0$, then
\[
R=R(A)
\geq
\min\{c_{0,1},c_*\}
>
R_*,
\]
which contradicts $R<R_*$. Hence
\[
U(0)<A_0<A_*.
\]
Since every positive radial shooting profile is strictly decreasing,
\[
\norm{U}_{L^\infty(B_R)}
=
U(0)
<
A_*.
\]
\textcolor{red}{Assuming, as in
Proposition~\ref{prop:smalluniqueness}, uniqueness in the class of
positive solutions under the bound
$\|U\|_{L^\infty(B_R)}<A_*$,}
there is at most one such positive \textcolor{red}{radial} solution.
Together with the existence of the small-amplitude radial solution,
this proves that problem~\eqref{eq:scaled} admits exactly one positive
radial solution $U_R$ for every $0<R<R_*$.
\end{proof}
    


%%%%%%%%%%%%%%%%%%%%%%%%%%%%%%%%%%%%%%%%%%%%%%%%%%%%%%%%%%%%%%%%%
  
%%%%%%%%%%%%%%%%%%%%%%%%%%%%%%%%%%%%%%%%%%%%%%%%%%%%%%%%%%%%%%%%%
\section{From radial uniqueness to full uniqueness}
%%%%%%%%%%%%%%%%%%%%%%%%%%%%%%%%%%%%%%%%%%%%%%%%%%%%%%%%%%%%%%%%%

The reaction
\[
f_\lambda(s)
=
\lambda(s^+)^q+(s^+)^p
\]
is not locally Lipschitz at $s=0$, since $0<q<1$. Nevertheless, it
admits the decomposition
\[
f_\lambda=f_1+f_2,
\qquad
f_1(s)=(s^+)^p,
\qquad
f_2(s)=\lambda(s^+)^q,
\]
where $f_1$ is locally Lipschitz and $f_2$ is nondecreasing.

This is precisely the structural decomposition treated by the moving-plane
symmetry theorem of Moreira dos Santos and Nornberg~\cite{SantosNornberg}:
the zero-order term is the sum of a locally Lipschitz function and a
nondecreasing function of the unknown. In the present scalar problem there are
no coupling or gradient terms, the operator is the uniformly elliptic
Laplacian, the reaction is independent of $x$ (hence in particular
nonincreasing in $|x|$), and the domain is the ball.
\textcolor{red}{Moreover, $f_\lambda$ is continuous on $\mathbb R$,
and the decomposition above satisfies the required structural assumptions
on the nonnegative range of the solution.}
Their theorem therefore
applies and shows that every positive classical solution of~\eqref{eq:ABC}
is radially symmetric about the origin and strictly decreasing in the radial
variable.

We can therefore complete the proof of the main result.

\begin{proof}[Proof of Theorem~\ref{thm:main}]
Choose $R_*>0$ as in Theorem~\ref{thm:radial} and set
\[
\lambda_*
:=
R_*^{\frac{2(p-q)}{p-1}}.
\]
If
\[
0<\lambda<\lambda_*,
\]
then, by~\eqref{eq:Rlambda},
\[
R
=
\lambda^{\frac{p-1}{2(p-q)}}
<
R_*.
\]
By Theorem~\ref{thm:radial}, the scaled
problem~\eqref{eq:scaled} has exactly one positive radial solution.
On the other hand, by the symmetry result quoted above, every
positive classical solution of~\eqref{eq:ABC} is radial.
\textcolor{red}{The same symmetry conclusion applies to the scaled
problem, since its nonlinearity has the same structural form.}
Since the
scaling~\eqref{eq:scaling} gives a one-to-one correspondence between
positive solutions of~\eqref{eq:ABC} and positive solutions
of~\eqref{eq:scaled}, it follows that~\eqref{eq:ABC} admits exactly
one positive classical solution, denoted by $u_\lambda$.

It remains to determine its size as $\lambda\to0^+$. Let
\[
A=U_R(0).
\]
\textcolor{red}{Since
\[
R=\lambda^{\frac{p-1}{2(p-q)}}\longrightarrow0
\qquad\text{as }\lambda\to0^+,
\]
Theorem~\ref{thm:radial} and the exclusion of amplitudes
$A\geq A_0$ in its proof imply that
\[
A=U_R(0)\longrightarrow0.
\]
Hence Proposition~\ref{prop:smallasymptotic} is applicable to this
family and gives}
\[
R
=
\rho_q A^{\frac{1-q}{2}}
\bigl(1+o(1)\bigr).
\]
Therefore
\[
A
=
\rho_q^{-\frac{2}{1-q}}
R^{\frac{2}{1-q}}
\bigl(1+o(1)\bigr).
\]
Using
\[
R
=
\lambda^{\frac{p-1}{2(p-q)}},
\]
we obtain
\[
A
=
\rho_q^{-\frac{2}{1-q}}
\lambda^{\frac{p-1}{(p-q)(1-q)}}
\bigl(1+o(1)\bigr).
\]
By the scaling~\eqref{eq:scaling},
\[
u_\lambda(0)
=
\lambda^{\frac1{p-q}}A.
\]
Hence
\begin{align*}
u_\lambda(0)
&=
\rho_q^{-\frac{2}{1-q}}
\lambda^{
\frac1{p-q}
+
\frac{p-1}{(p-q)(1-q)}
}
\bigl(1+o(1)\bigr)
\\
&=
\rho_q^{-\frac{2}{1-q}}
\lambda^{\frac1{1-q}}
\bigl(1+o(1)\bigr),
\end{align*}
\textcolor{red}{where we used
\[
\frac1{p-q}
+
\frac{p-1}{(p-q)(1-q)}
=
\frac{(1-q)+(p-1)}{(p-q)(1-q)}
=
\frac1{1-q}.
\]}
Since $u_\lambda$ is radial and strictly decreasing,
\[
\norm{u_\lambda}_{L^\infty(B_1)}
=
u_\lambda(0).
\]
Consequently,
\begin{equation}
\label{eq:main-asymptotic}
\norm{u_\lambda}_{L^\infty(B_1)}
\sim
\rho_q^{-\frac{2}{1-q}}
\lambda^{\frac1{1-q}}
\qquad
\text{as }\lambda\to0^+.
\end{equation}
This completes the proof.
\end{proof}
    


%%%%%%%%%%%%%%%%%%%%%%%%%%%%%%%%%%%%%%%%%%%%%%%%%%%%%%%%%%%%%%%%%
  
%%%%%%%%%%%%%%%%%%%%%%%%%%%%%%%%%%%%%%%%%%%%%%%%%%%%%%%%%%%%%%%%%
\section{Comments on the mechanism}
%%%%%%%%%%%%%%%%%%%%%%%%%%%%%%%%%%%%%%%%%%%%%%%%%%%%%%%%%%%%%%%%%
\begin{remark}
At the Emden--Fowler level, the damping coefficient is
\[
A
=
N-2-\frac4{p-1}.
\]
Thus
\[
A>0
\quad\Longleftrightarrow\quad
p>\frac{N+2}{N-2}.
\]
In the supercritical regime, we also have $B>0$, and the
linearization at the positive Lane--Emden equilibrium $(\ell,0)$
has characteristic polynomial
\[
\mu^2+A\mu+(p-1)B.
\]
Since
\[
A>0
\qquad\text{and}\qquad
(p-1)B>0,
\]
both characteristic roots have strictly negative real parts.
Thus $(\ell,0)$ is exponentially asymptotically stable, and this
stability is the key ingredient in the proof of
Lemma~\ref{lem:S-lower}.

At the critical exponent,
\[
p=\frac{N+2}{N-2},
\]
one has $A=0$, so the positive equilibrium is no longer
asymptotically stable by this damping mechanism.
\textcolor{red}{Indeed, the linearized equation becomes
\[
z''+(p-1)Bz=0,
\]
whose characteristic roots are
\[
\mu_\pm=\pm i\sqrt{(p-1)B}.
\]
Moreover, the Emden--Fowler equation itself becomes the conservative
equation
\[
x''-Bx+x^p=0,
\]
for which the energy
\[
E(x,x')
=
\frac12(x')^2-\frac{B}{2}x^2+\frac1{p+1}x^{p+1}
\]
is constant along trajectories.}
Consequently, the
large-amplitude argument used here breaks down.
\textcolor{red}{More precisely, the robust trapping mechanism of
Lemma~\ref{lem:robust}, which relies on exponential asymptotic
stability of $(\ell,0)$, is no longer available when $A=0$.
Thus the critical case requires a different analysis and should not
be inferred from the present supercritical argument. This distinction
is also compatible with the qualitatively different behavior known
for critical concave--convex problems of
Ambrosetti--Brezis--Cerami type.}
\end{remark}

\begin{remark}
The estimate
\[
S(\varepsilon)
\geq
c\varepsilon^{-\frac{p-1}{2(p-q)}}
\]
has precisely the scale needed to prevent the original first zero
$R(A)$ from collapsing to zero as $A\to\infty$. Indeed, the
large-amplitude rescaling gives
\(
R(A)
=
A^{-\frac{p-1}{2}}
S(A^{q-p}),
\)
and therefore
\begin{align*}
R(A)
&\geq
cA^{-\frac{p-1}{2}}
\left(A^{q-p}\right)^{-\frac{p-1}{2(p-q)}}
\\
&=
\textcolor{red}{
cA^{-\frac{p-1}{2}}
A^{-(q-p)\frac{p-1}{2(p-q)}}
}
\\
&=
cA^{-\frac{p-1}{2}}
A^{\frac{p-1}{2}}
=
c,
\end{align*}
\textcolor{red}{where we used $q-p=-(p-q)$.}
Thus the two scaling factors cancel exactly, which is the mechanism
behind the uniform lower bound for $R(A)$ at large amplitudes.
\end{remark}
    


%%%%%%%%%%%%%%%%%%%%%%%%%%%%%%%%%%%%%%%%%%%%%%%%%%%%%%%%%%%%%%%%%


%%%%%%%%%%%%%%%%%%%%%%%%%%%%%%%%%%%%%%%%%%%%%%%%%%%%%%%%%%%%%%%%%
\section*{Acknowledgment of scope}
%%%%%%%%%%%%%%%%%%%%%%%%%%%%%%%%%%%%%%%%%%%%%%%%%%%%%%%%%%%%%%%%%

The uniqueness theorem proved above concerns positive solutions in a
ball. The proof uses both the radial shooting geometry of the ball and
a radial symmetry theorem for positive solutions. No corresponding
uniqueness statement is asserted here for general bounded domains.

%%%%%%%%%%%%%%%%%%%%%%%%%%%%%%%%%%%%%%%%%%%%%%%%%%%%%%%%%%%%%%%%%
\begin{thebibliography}{99}

\bibitem{ABC}
A.~Ambrosetti, H.~Brezis and G.~Cerami,
\newblock Combined effects of concave and convex nonlinearities in some elliptic problems,
\newblock \emph{J. Funct. Anal.} \textbf{122} (1994), 519--543.

\bibitem{GNN}
B.~Gidas, W.-M.~Ni and L.~Nirenberg,
\newblock Symmetry and related properties via the maximum principle,
\newblock \emph{Comm. Math. Phys.} \textbf{68} (1979), 209--243.

\bibitem{Ma1996}
R.~Ma,
\newblock On a conjecture concerning the multiplicity of positive solutions of elliptic problems,
\newblock \emph{Nonlinear Anal.} \textbf{27} (1996), 775--780.

\bibitem{NiNussbaum}
W.-M.~Ni and R.~D.~Nussbaum,
\newblock Uniqueness and nonuniqueness for positive radial solutions of
$\Delta u+f(u,r)=0$,
\newblock \emph{Comm. Pure Appl. Math.} \textbf{38} (1985), 67--108.

\bibitem{SantosNornberg}
E.~Moreira dos Santos and G.~Nornberg,
\newblock Symmetry properties of positive solutions for fully nonlinear elliptic systems,
\newblock \emph{J. Differential Equations} \textbf{269} (2020), 4175--4191.
\textcolor{red}{
\bibitem{ZhaoZhong2002}
P.~Zhao and C.~Zhong,
\newblock Positive solutions of elliptic equations involving both
supercritical and sublinear growth,
\newblock \emph{Houston J. Math.} \textbf{28} (2002), 649--663.
}

\end{thebibliography}

\end{document}